\subsection{The generic upper bound from a finite slice}
\label{subsec:selmer-upper-multiplicity}

Keep the algebra, equations, and three Rees modules of
Proposition~\ref{prop:selmer-admissibility}. Retain the finite
degree-zero quotient, negative cotangent spanning, and fixed-point
separation assumptions. Assume also that the tangent space at every
two-sided $K$-point of $A/Q$ has dimension at most one. Set
\[
 Q_0=Q\cap B,\qquad \mathfrak d=A_{-2}A_2,\qquad q=r+1+r'\ge1,
 \qquad Q_{\mathrm{sq}}=(F_i^2,F_*^2,G_j,E_k^2,E'_l,Z_i).
\]
Let $J\subset\mathcal R$ be homogeneous, containing the squared
sections $\boldsymbol s$ and the $Z_i$, and write
$M_{\boldsymbol d}=(\mathcal R/J)_{\boldsymbol d}$.

\begin{proposition}[Generic support and multiplicity]
\label{prop:selmer-generic-multiplicity}
In all sufficiently large multidegrees, $M_{\boldsymbol d}$ and
$B/\operatorname{Ann}_B M_{\boldsymbol d}$ are finite over
$\mathcal O$. If $P\cap\mathcal O=0$ and
$(M_{\boldsymbol d})_P\ne0$, then $Q_0\subset P$,
$\mathfrak d\not\subset P$, and
\[
 \ell_{B_P}\bigl((M_{\boldsymbol d})_P\bigr)\le2^q.
\]
Every embedding $\kappa(P)\hookrightarrow K$ is the restriction of
a two-sided point of $A/Q$.
\end{proposition}

\begin{proof}
The clearing identities in the admissibility proof show that every
relevant prime over $J$ contains $Q_0\mathcal R$. Thus
$Q_0\mathcal R I_{\mathrm{irr}}\subset\sqrt J$, and Noetherianity
gives $(Q_0\mathcal R I_{\mathrm{irr}})^n\subset J$ for some $n$.
Standardness gives
$M_{\boldsymbol d}=(I_{\mathrm{irr}}^nM)_{\boldsymbol d}$ in large
degrees, hence $Q_0^nM_{\boldsymbol d}=0$.
The filtration of $B/Q_0^n$ by powers of $Q_0$ has finite
$B/Q_0$-modules as successive quotients, so it is finite over
$\mathcal O$. This proves both finiteness assertions.

Apply the same prime argument to $J+\mathfrak d\mathcal R$.
Every relevant prime contracts to $\mathfrak m_B$: a
characteristic-zero contraction killing $\mathfrak d$ is excluded by
the half-point construction and the cotangent argument. Therefore
\[
 (\mathfrak m_B\mathcal R I_{\mathrm{irr}})^{n'}
 \subset J+\mathfrak d\mathcal R
\]
for some $n'$. In a common eventual range,
$\mathfrak m_B^{n'}M_{\boldsymbol d}\subset\mathfrak dM_{\boldsymbol d}$.
At a prime $P$ containing $\mathfrak d$ with $P\cap\mathcal O=0$,
a uniformizer is invertible. Localization and Nakayama's lemma give
$(M_{\boldsymbol d})_P=0$. For every remaining generic prime and
embedding $\sigma:\kappa(P)\hookrightarrow K$, the nonzero
finite-type graded algebra $(A/Q)\otimes_{B,\sigma}K$ has
degree-zero part $K$. A $K$-point exists, and it is two-sided because
$\sigma(\mathfrak d)\ne0$.

Fix such a prime and choose $f\in A_{-2},g\in A_2$ with
$fg\notin P$. Both elements are units in $A_P$. Put
\[
 D=(B/Q_0)_P,\qquad C=(A/Q)_P,\qquad S=C/(f-1).
\]
The ring $D$ is local Artinian and finite over $F$. The graded pieces
of $C$ satisfy $C_{2j}=f^{-j}D$ and $C_{2j+1}=f^{-j}C_1$.
Taking the direct sum of these Laurent powers and imposing $f=1$
gives the actual module decomposition $S=D\oplus C_1$, so $D$
embeds in $S$. The ring $S$ is finite over $D$: after normalizing the
finitely many generators by powers of $f$, each even generator lies
in $D$ and each odd generator has its square in $D$.
The parity involution has invariant subring exactly $D$.

At a $K$-point of $S\otimes_F K$, the Euler derivation
$c_j\mapsto j\xi(c_j)$ is a tangent vector of $A/Q$ taking $f$
to $-2$. Every slice tangent takes $f$ to zero, and the tangent
dimension assumption therefore makes it zero. Each local maximal
ideal of the finite algebra $S\otimes_F K$ has
$\mathfrak m/\mathfrak m^2=0$, hence is zero by Nakayama.
Thus $S\otimes_F K$ is reduced. Its subring $D\otimes_F K$ is
reduced as well, so $D=\kappa(P)$ is a field.

For every embedding $D\hookrightarrow K$, the finite reduced
algebra $S\otimes_D K$ has invariant subring $K$. The involution
permutes its geometric points, and its invariant algebra has one
factor per orbit. There is exactly one orbit, of size one or two.
The even and odd dimensions are consequently
\begin{equation}
 (m_0,m_1)=(1,0)\quad\hbox{or}\quad(1,1),\qquad
 \ell_{B_P}\bigl((A/Q)_{j,P}\bigr)=m_{j\bmod2}.
 \label{eq:selmer-slice-parities}
\end{equation}
All these arguments take place in characteristic zero. They impose
no restriction on the residue characteristic of $\mathcal O$.

In $A_P/Q_{\mathrm{sq}}$, the ideal $U$ generated by the $q$ selected
odd equations satisfies $U^{q+1}=0$. Each quotient $U^k/U^{k+1}$
is generated over $A_P/Q$ by the squarefree products of $k$ selected
equations. This gives surjections from the corresponding shifted
copies of $A_P/Q$; linear independence is unnecessary. The filtration
by powers of $U$ gives, in every weight,
\begin{equation}
 \ell_{B_P}\bigl((A/Q_{\mathrm{sq}})_{w,P}\bigr)
 \le\sum_{I\subset\{1,\ldots,q\}}m_{(w-|I|)\bmod2}
 =2^{q-1}(m_0+m_1)\le2^q.
 \label{eq:selmer-squarefree-multiplicity}
\end{equation}

It remains to compare the actual Rees piece with this weight piece.
The same units give
$(N_0)_P=fB_P$, $(N_1)_P=gB_P$, and $(N_2)_P=fB_P$;
the last equality uses $f^2g=(fg)f\in N_2$. Therefore
\[
 \mathcal R_P=B_P[X_0,X_1,X_2],\qquad
 X_0=ft_0,\quad X_1=gt_1,\quad X_2=ft_2.
\]
Let $H\subset B_P$ be generated by
$F_i^2/f,G_j/f,E_k^2/g,E'_l/g,F_*^2/f,Z_i$.
When all $d_i\ge1$, the degree-$\boldsymbol d$ quotient by the
specified sections and $Z_i$ is $B_P/H$, since every equation is
available in that degree. The further ideal $J$ gives a quotient.
On the other hand $Q_{\mathrm{sq}}A_P=HA_P$, and the even weight
$w=2(d_1-d_0-d_2)$ has unit generator $f^{d_0+d_2}g^{d_1}$.
Thus its quotient piece is also $B_P/H$, giving
\[
 (A/Q_{\mathrm{sq}})_{w,P}\twoheadrightarrow
 (M_{\boldsymbol d})_P\qquad(d_i\ge1).
\]
Together with \eqref{eq:selmer-squarefree-multiplicity}, this proves
the bound. No comparison before localization or on a missing
coordinate face is required.
\end{proof}

At a single-point slice, \eqref{eq:selmer-squarefree-multiplicity}
improves the coefficient to $2^{q-1}$. Which slice types arise in the
arithmetic model remains a separate question.

\paragraph{Completion of the abstract numerical comparison.}
Let $a\in B$ have augmentation zero and be nonzero at every
two-sided equation point. Let $\Sigma$ be the distinct restrictions
to $B$ of a finite set of torus-orbit representatives covering those
points. Proposition~\ref{prop:selmer-generic-multiplicity} and the
determinant--Euler identity give
\[
 \chi_B(a;M_{\boldsymbol d})
 =\sum_{\sigma\in\Sigma}m_{\sigma,\boldsymbol d}v_\pi(\sigma(a))
 \le 2^{r+1+r'}\sum_{\sigma\in\Sigma}v_\pi(\sigma(a))
\]
eventually. The valuations are nonnegative because $B/Q_0$ is
finite over $\mathcal O$. Retain all the chart and Smith hypotheses
of the lower comparison, including $r\ge1$, the prescribed chart
weights, vanishing on the $Z_i$, and the separate condition
$\varphi(a)\ne0$. The preceding results then prove
\begin{equation}
 \sum_i c_i\le2\sum_{\sigma\in\Sigma}v_\pi(\sigma(a)).
 \label{eq:selmer-complete-abstract-core}
\end{equation}
The coefficient on an individual valuation improves to one when its
slice has one point. Replacing distinct characters by the original
orbit representatives only increases the nonnegative sum. The next
theorem derives the required chart normalization and finite orbit cover
from the original model hypotheses. Point nonvanishing alone does not
establish formal-chart nonvanishing for the original element; an
equation-valued adjustment supplies a suitable representative.
